\chapter{Theory and Fundamentals}
This chapter covers the fundamental concepts and technologies used in the FoodSense AI project, including sensors, machine learning architectures, and the tools and software employed.

\section{Sensors and Hardware}
The FoodSense AI system uses two main sensors and an Arduino microcontroller for data acquisition.

\subsection{MQ135 Gas Sensor}
The MQ135 is a low-cost metal-oxide gas sensor sensitive to ammonia, sulfide, benzene vapor, and other harmful gases. It operates by measuring the change in resistance of a tin dioxide (SnO₂) layer when exposed to target gases. The sensor requires a load resistor (typically 10 kΩ) and operates at 5V DC.

The resistance of the MQ135 sensor decreases as the concentration of target gases increases. The relationship between sensor resistance and gas concentration follows a power-law:
\[
C = A \times \left( \frac{R_s}{R_0} \right)^B
\]
where \(C\) is the gas concentration, \(R_s\) is the sensor resistance, \(R_0\) is the sensor resistance in clean air, and \(A\) and \(B\) are calibration constants.

\subsection{DHT11 Temperature and Humidity Sensor}
The DHT11 is a basic, ultra-low-cost digital temperature and humidity sensor. It uses a capacitive humidity sensor and a thermistor to measure the surrounding air, and outputs a digital signal on the data pin. The sensor provides temperature readings in the range of 0-50°C with ±2°C accuracy, and humidity readings in the range of 20-90\% RH with ±5\% accuracy.

\subsection{Arduino Microcontroller}
Arduino is an open-source electronics platform based on easy-to-use hardware and software. The FoodSense AI system uses an Arduino Uno, which is based on the ATmega328P microcontroller. The Arduino reads data from the MQ135 and DHT11 sensors, performs on-device compensation, and transmits the data via serial communication in CSV format.

\section{Machine Learning Architectures}
The FoodSense AI system uses two main machine learning models: a Multi-Layer Perceptron (MLP) for classification and a Long Short-Term Memory (LSTM) network with temporal attention for regression.

\subsection{Multi-Layer Perceptron (MLP)}
An MLP is a class of feedforward artificial neural network that consists of at least three layers of nodes: an input layer, a hidden layer, and an output layer. Except for the input nodes, each node is a neuron that uses a nonlinear activation function. MLPs are widely used for classification and regression tasks.

The MLP used in this project has the following architecture:
\begin{itemize}
    \item Input layer: 128 features (from the UCI Gas Sensor Array Drift Dataset)
    \item Hidden layer 1: Linear(128) + BatchNorm + ReLU + Dropout(0.3)
    \item Hidden layer 2: Linear(64) + BatchNorm + ReLU + Dropout(0.2)
    \item Hidden layer 3: Linear(32) + ReLU
    \item Output layer: Linear(1) → BCEWithLogitsLoss
\end{itemize}

\subsection{Long Short-Term Memory (LSTM)}
LSTMs are a type of recurrent neural network (RNN) architecture used in the field of deep learning. Unlike standard feedforward neural networks, LSTMs have feedback connections, making them suitable for processing sequences of data. LSTMs are well-suited for time-series prediction tasks because they can learn long-term dependencies.

The LSTM used in this project includes a temporal attention mechanism, which allows the model to weight the importance of different timesteps in the input sequence. This is particularly useful for food spoilage detection, as not all past sensor readings are equally predictive of future spoilage.

The LSTM architecture used is:
\begin{itemize}
    \item Input sequence: last 30 readings (1 hour context window)
    \item LSTM layer: hidden size 128, 2 layers, dropout 0.3
    \item Temporal attention layer
    \item Linear layer: Linear(64) + ReLU + Dropout(0.2)
    \item Output layer: Linear(1) + Sigmoid
\end{itemize}

\section{Tools and Software}
Several tools and software packages were used in the development of FoodSense AI.

\subsection{Python and PyTorch}
Python is a high-level, general-purpose programming language widely used in machine learning and data science. PyTorch is an open-source machine learning library for Python, developed by Meta AI. PyTorch provides tensor computation with strong GPU acceleration and deep neural networks built on a tape-based autograd system.

\subsection{Arduino IDE}
The Arduino Integrated Development Environment (IDE) is a cross-platform application written in Java, which is used to write and upload programs to Arduino-compatible boards. The Arduino IDE supports C and C++ programming languages.

\subsection{HTML, CSS, and JavaScript}
HTML, CSS, and JavaScript are the core technologies of the World Wide Web. HTML provides the structure of web pages, CSS describes the presentation of web pages, and JavaScript adds interactivity to web pages. The FoodSense AI dashboard is built using these technologies, with Tailwind CSS for styling.

\section{Summary}
This chapter covered the fundamental concepts and technologies used in the FoodSense AI project, including the MQ135 and DHT11 sensors, Arduino microcontroller, MLP and LSTM machine learning architectures, and the tools and software employed. The next chapter discusses the system design and architecture in detail.
